\section{Managerial Process}

Our three-member team follows a few simple rules for priorities, communication and change \cite{spp1}.

\subsection{Objectives and Priorities}

Our objective is a system that lets a teacher run a fair lab exam from start to finish. When time is short, we protect work in this order:
\begin{enumerate}
    \item \textbf{Exam integrity:} device binding, eligibility and seat checks, and lockdown (M1, M9, M4).
    \item \textbf{No lost work:} autosave, encrypted backup and server deadline enforcement (M6).
    \item \textbf{Live supervision:} monitoring client and proctoring dashboard (M5, M10).
    \item \textbf{Evaluation support:} test-case grading, theory scoring and plagiarism checks (M11).
    \item \textbf{Reporting:} audit logs, analytics and PDF reports (M12).
\end{enumerate}
Lower items are cut first.

\subsection{Assumptions and Constraints}

\textbf{Assumptions}
\begin{itemize}
    \item Lab PCs run Windows 10 or 11, and we can install the client with administrator rights.
    \item Each PC reaches the API over HTTPS; short outages may happen.
    \item Each student uses one assigned workstation, with a webcam when the exam needs one.
    \item Teachers supply test cases and model answers for automated grading.
\end{itemize}

\textbf{Constraints}
\begin{itemize}
    \item Three members and 32 weeks over two semesters.
    \item No budget for paid services: we use the Neon free tier, and AI features run offline on the CPU (TF-IDF scoring, OpenCV YuNet face detection). LLM grading is off by default.
    \item The client is Windows-only (WPF) and cannot block Ctrl+Alt+Del.
    \item Code execution needs Docker on the grading server.
    \item For privacy, there is no face recognition and no stored biometric data.
\end{itemize}

\subsection{Communication, Review and Change Control}

We meet at the start of each sprint to set goals and at the end to review the results. Between meetings we post short written check-ins \cite{scrum1}, and a blocker unsolved after 24 hours goes to a team discussion. Supervisor feedback goes into the next sprint's backlog. A proposed change is accepted only if it improves the project, fits in the sprint and does not put a working module at risk. Otherwise it goes to the future-work list.
